% paper/sections/04e_fccd.tex
%
% §4.5  FCCD：面ジェット共通プリミティブ（HFE / GFM / 分相PPE）
%        4 設計原理・4 次導出および行列定式化を統合．
%        H-01 歴史的診断は付録 §B.4（appendix_ccd_h01.tex）に退避．
%
% ═══════════════════════════════════════════════════════════════════════════════
\subsection{\texorpdfstring{{\FCCD}：面ジェット共通プリミティブ}{FCCD：面ジェット共通プリミティブ}}
\label{sec:fccd_def}
% ═══════════════════════════════════════════════════════════════════════════════

本節では\textbf{面中心 Combined Compact Difference}（{\FCCD}）演算子族を，
4 設計原理，4 次 cancellation 導出，行列定式化，周期 DFT 解析を順に提示して定義する．
{\FCCD} は単独スキームではなく，複数の界面消費点で共有される
\textbf{共通面ジェットプリミティブ} $\mathcal{J}_f(u)=(u_f, u'_f, u''_f)$ を
$\Ord{H^4}$ で供給する契約（shared contract）として本稿では位置付けられる．

\noindent\textbf{本稿での 4 つの再利用先（R1)-(R4)．}
\begin{enumerate}[label=\textbf{(R\arabic*)}]
  \item HFE 上流点再構築（§~\ref{sec:advection_motivation}）における面ジェット消費；
  \item GFM ghost pressure jet（§~\ref{sec:split_ppe_framework}）での $[p]=\sigma\kappa$ および $[\beta_f p'_f]=0$（$\beta_f=(1/\rho)_f$）条件課嵌；
  \item 分相 PPE 楕円作用素 $G_h^{bf}$ と $D_h^{bf} = -(G_h^{bf})^{\!*}$（§~\ref{sec:split_ppe_framework}）の離散随伴対；
  \item 粘性 stress-divergence の界面バンド（§~\ref{sec:viscous_3layer}）における面法線・接線高次微分．
\end{enumerate}
バルク演算子は節点 CCD の $\Ord{h^6}$ を維持し，
{\FCCD} の $\Ord{H^4}$ は上記 4 界面消費点でのみ展開される．
この意味で {\FCCD} は bulk 精度の縮退ではなく，
界面近傍のメトリック整合を担保する\textbf{設計選択}である．

\noindent\textbf{歴史的経緯の保全．}
{\FCCD} 初版の動機は FVM 面平均勾配 $\mathcal{G}^{\text{adj}}$ との
メトリック整合（H-01 診断）にあったが，本稿の純 FCCD 設計では FVM が除外される．
歴史的診断 H-01 および Exp-1/Exp-2 実測値は付録~\ref{sec:fccd_motivation_h01}
（§B.4 H-01 診断の歴史的ケーススタディ）に保全されている．

% -------------------------------------------------------------------------------
\subsubsection{4 つの設計原理}
\label{sec:fccd_design_principles}
% -------------------------------------------------------------------------------

{\FCCD} は以下の 4 原理に従う：

\begin{enumerate}[label=\textbf{P-F\arabic*.}]
\item \textbf{風上限定．}
      $a>0$ において $\{u_{i-1}, u_i\}$ のみ使用，
      中心・下流成分に依存しない．
\item \textbf{ステンシル非拡大．}
      $u_{i-2}$ 以遠は参照せず，3 点コンパクト性を CCD から継承．
\item \textbf{Combined Compact 構造．}
      高階補助未知数（$u'''_f$）を導入し，
      Chu--Fan 同様の\textbf{微分関係式の組合せ}で精度を取得．
\item \textbf{面中心化．}
      一次量は全て面 $f = i-1/2$ で定義，
      HFE 上流点・GFM ジャンプ行・分相 PPE 楕円演算子と\textbf{同一面ロカス}．
\end{enumerate}

% -------------------------------------------------------------------------------
\subsubsection{4 次 Cancellation 導出}
\label{sec:fccd_derivation}
% -------------------------------------------------------------------------------

1 次差分の面上展開：
\begin{equation}
  \frac{u_i - u_{i-1}}{\Delta x}
  \;=\; u'_f + \frac{\Delta x^2}{24}\,u'''_f + \Ord{\Delta x^4}.
  \label{eq:fccd_upwind_expansion}
\end{equation}
Chu--Fan 流の組合せ演算子を定義：
\begin{equation}
  D^{\FCCD} u_f
  \;=\; \frac{u_i - u_{i-1}}{\Delta x}
  \;-\; \lambda\,\Delta x^2\,u'''_f.
  \label{eq:fccd_operator}
\end{equation}
式~\eqref{eq:fccd_upwind_expansion} を代入して \textbf{cancellation 条件}
$\lambda = 1/24$ を課せば：
\begin{equation}
  \boxed{\;
  D^{\FCCD} u_f \;=\; u'_f + \Ord{\Delta x^4}
  \;}
  \label{eq:fccd_fourth_order}
\end{equation}
となり，ステンシル拡大・風上構造の放棄なしに\textbf{4 次精度}を獲得する．
$u'''_f$ は近傍節点の $q_i = u''_i$（CCD 出力）から
$u'''_f \approx (q_i - q_{i-1})/\Delta x$ の 2 次中心差で閉じる．

% -------------------------------------------------------------------------------
\subsubsection{行列定式化}
\label{sec:fccd_matrix}
% -------------------------------------------------------------------------------

行列表記を用いると，面勾配ベクトル
$\mathbf{d}^{\FCCD} \in \mathbb{R}^N$ は節点値 $\mathbf{u}\in\mathbb{R}^{N+1}$ から：
\begin{equation}
  \mathbf{d}^{\FCCD}
  \;=\; \mathbf{D}_1\,\mathbf{u} - \mathbf{D}_2\,\mathbf{q},
  \qquad \mathbf{q} = \mathbf{S}_{\CCD}\,\mathbf{u},
  \label{eq:fccd_matrix_system}
\end{equation}
で得られる．ここで
$\mathbf{D}_1, \mathbf{D}_2 \in \mathbb{R}^{N\times(N+1)}$ は面↔節点
\textbf{bidiagonal}（行あたり 2 非零）：
\begin{align}
  (\mathbf{D}_1)_{f_{i-1/2},\,i-1} &= -1/H, &\quad (\mathbf{D}_1)_{f_{i-1/2},\,i} &= +1/H, \label{eq:fccd_D1}\\
  (\mathbf{D}_2)_{f_{i-1/2},\,i-1} &= -H/24, &\quad (\mathbf{D}_2)_{f_{i-1/2},\,i} &= +H/24. \label{eq:fccd_D2}
\end{align}
$\mathbf{S}_{\CCD} = \Pi_q\,\mathbf{A}_{\CCD}^{-1}\mathbf{B}_{\CCD}$ は
Chu--Fan 結合システムを節点 2 次導関数へ射影した作用素で，
\textbf{形式上は密行列}だが実行時は block-Thomas（壁）または
block-circulant LU（周期）で $\Ord{N}$ 計算．
従って合成演算子：
\begin{equation}
  \boxed{\;
  \mathbf{M}^{\FCCD}
  \;=\; \mathbf{D}_1 - \mathbf{D}_2\,\mathbf{S}_{\CCD}
  \;\in\; \mathbb{R}^{N\times(N+1)}
  \;}
  \label{eq:fccd_composite_matrix}
\end{equation}
の適用は 2 段階：
(i) $\mathbf{q} \leftarrow \mathbf{S}_{\CCD}\mathbf{u}$（既存 CCD ソルバ再利用），
(ii) $\mathbf{d}^{\FCCD} \leftarrow \mathbf{D}_1\mathbf{u} - \mathbf{D}_2\mathbf{q}$（面上局所ステンシル）．
全体コスト $\Ord{N}$，追加メモリは $\mathbf{q}$ のみ．

% -------------------------------------------------------------------------------
\subsubsection{周期境界での DFT 解析}
\label{sec:fccd_periodic_dft}
% -------------------------------------------------------------------------------

周期格子で $\mathbf{M}^{\FCCD}$ は\textbf{ブロック循環行列}（block-circulant）となり，
離散 Fourier 変換で対角化される．
波数 $\xi = k H$ に対する Fourier シンボル：
\begin{equation}
  \hat M^{\FCCD}(\xi)
  \;=\; \frac{2i\sin(\xi/2)}{H}\cdot e^{-i\xi/2}
        \;\cdot\;
        \left[1 - \frac{\xi^2}{24}\,\hat S_{\CCD}(\xi)\right]
  \label{eq:fccd_fourier_symbol}
\end{equation}
（$\hat S_{\CCD}(\xi) = -\omega_2(\xi)^2/H^2$）
を Taylor 展開して $\hat M^{\FCCD}(\xi) = i\xi/H + \Ord{\xi^5/H}$ を得，
式~\eqref{eq:fccd_fourth_order} の 4 次精度が
Fourier レベルでも確認される．

% -------------------------------------------------------------------------------
\subsubsection{非一様格子への拡張（概要）}
\label{sec:fccd_nonuniform_sketch}
% -------------------------------------------------------------------------------

非一様格子（$H_i$ 可変，面位置パラメータ $\theta_i$）での
cancellation 係数 $\lambda_i$ と混合項係数 $\mu_i$ は格子幾何に依存：
\begin{equation}
  \mu_i = \theta_i - \tfrac{1}{2},
  \qquad
  \lambda_i = \frac{1 - 3\theta_i(1-\theta_i)}{6}
  \label{eq:fccd_nonuniform_coeffs}
\end{equation}
．合成演算子は
$\mathbf{M}^{\FCCD,\text{nu}} = \mathbf{D}_1^{(H)} - (\mathbf{D}_\mu^{(H\theta)} + \mathbf{D}_\lambda^{(H)})\,\mathbf{S}_{\CCD}$
で，一様極限 $\theta_i\equiv 1/2 \Rightarrow \mu_i\equiv 0, \lambda_i \equiv 1/24$ で
式~\eqref{eq:fccd_composite_matrix} を回復する．
完全な係数行列と切断誤差解析は §\ref{sec:grid}（非一様格子）で扱う．

% -------------------------------------------------------------------------------
\subsubsection{Balanced-Force 補正（面ジェット共通メトリック）}
\label{sec:fccd_h01_remedy}
% -------------------------------------------------------------------------------

{\FCCD} を圧力勾配と界面力勾配の\textbf{両方}に適用：
\begin{equation}
  \mathcal{G}\,p \;\equiv\; D^{\FCCD}\,p,
  \qquad
  \bnabla\psi \;\equiv\; D^{\FCCD}\,\psi,
  \label{eq:fccd_bf_unification}
\end{equation}
すると両演算子は同一面ロカス・同一メトリック・同一 Fourier シンボルとなり
Balanced-Force 残差は：
\begin{equation}
  |\mathrm{BF}_{\text{res}}|
  \;=\; \bigl\|\mathcal{G}\,p_{\text{YL}} - \sigma\kappa\,\bnabla\psi\bigr\|_\infty
  \;=\; \Ord{\Delta x^4}
  \label{eq:fccd_bf_residual_order}
\end{equation}
まで圧縮される．
4 役割 (R1)-(R4) が同一面ジェット $\mathcal{J}_f$ を共有することで，
圧力勾配・界面力・HFE 上流点・GFM ジャンプ条件の相互整合が
$\Ord{\Delta x^4}$ レベルで担保される
（初版の H-01 救済動機は付録~\ref{sec:fccd_motivation_h01} を参照）．

より運用的な Balanced-Force 設計原則（7 原理 P-1..P-7）は
§\ref{sec:balanced_force} で統一する；
{\FCCD} は P-1（演算子ロカス整合）・P-3（$D_h, G_h$ 離散随伴）・
P-5（フィルタ抑制下での安定性）の 3 原理を同時に満たす
\textbf{最小構成}である（証明は §\ref{sec:balanced_force} 参照）．

% -------------------------------------------------------------------------------
\subsubsection{導入の前提と実装チェック}
\label{sec:fccd_poc_requirements}
% -------------------------------------------------------------------------------

{\FCCD} 採用にあたり必要となる 3 点は以下のとおり：

\begin{itemize}
  \item \textbf{非一様一般化}（§\ref{sec:fccd_nonuniform_sketch}）は
        §\ref{sec:fccd_nonuniform} および ridge-Eikonal 非一様拡張
        （§\ref{sec:ridge_eikonal_nu_grid}）で完了．
  \item \textbf{壁境界}は Option III（Neumann）・Option IV（Dirichlet）を
        §\ref{sec:ccd_bc} で扱う．
  \item \textbf{擬時間 PPE 整合}は §\ref{sec:pressure} で
        純 {\FCCD} PPE として再構成する．
\end{itemize}

% ═══════════════════════════════════════════════════════════════════════════════
